
\section{Convex Biproduct Categories}\label{sec:cbproducts}  
In this section we introduce categories with convex biproducts, study the underlying monoidal algebra and illustrate a result which can almost be thought of as an analogue of Fox's theorem.


\subsection{Convex products}\label{ssec:convexproducts} As is the case for all Kleisli categories of monads on $\Sets$, the category $\KlD$ inherits coproducts from $\Sets$: the operation $\oplus$  in \eqref{ex:products} serves as a coproduct in $\KlD$. Our initial observation is that $\oplus$ satisfies an additional universal property, which we refer to as the \emph{convex product}, described below.

%As with all Kleisli categories of monads on $\Sets$, the category $\KlD$ inherits coproducts from $\Sets$: the operation $\oplus$, defined as in \eqref{ex:products}, serves as a coproduct in $\KlD$. Our key observation is that $\oplus$ also satisfies an additional universal property, which we describe below.
%So far, we have recalled that $\KlD$ is $\Cat{PCA}$-enriched. Like all Kleisli category of a monad over $\Sets$, $\KlD$ inherits coproducts from $\Sets$ \cite{}. In this section, we observe that such coproducts in $\KlD$ carries what we named \emph{convex products}.

\begin{definition}\label{def:convprod}
Let $X_1$, $X_2$  be two objects of a $\Cat{PCA}$-enriched category $\Cat{C}$. The \emph{convex product} of $X_1$ and $X_2$ is an object $Z$ with two arrows $\pi_1\colon Z\to X_1$ and $\pi_2\colon Z\to X_2$ satisfying the following property: for all $p_1,p_2 \in [0,1]$ such that $p_1+p_2\le 1$ and all arrows $f_1\colon A\to X_1$, $f_2\colon A \to X_2$, there exists a unique arrow $h \colon A \to Z$ %such that $h;\pi_1 = p_1\cdot f$ and $h;\pi_2 = p_2 \cdot g$. 
making the following diagram commute: %\marginpar{Cipriano, puoi meettere l'etichetta al centro della freccia?}
\[\begin{tikzcd}
	& A \\
	{X_1} & Z & {X_2}
	\arrow["{p_1\cdot f_1}"', from=1-2, to=2-1]
	\arrow["h", dashed, from=1-2, to=2-2]
	\arrow["{p_2\cdot f_2}", from=1-2, to=2-3]
	\arrow["{\pi_1}", from=2-2, to=2-1]
	\arrow["{\pi_2}"', from=2-2, to=2-3]
\end{tikzcd}\]
\end{definition}
%\marginpar{Pippo: Mi accorgo adesso che per $p=0$ la corrispondenza con i prodotti veri non vale...}

Similarly, the convex product of $n$ objects $X_1,\ldots,X_n$ is an object $Z$ with arrows $\pi_i\colon Z\to X_i$ for $i=1,\ldots,n$ satisfying the following property: for all $p_1, \dots, p_n \in [0,1]$ where $\sum_{i=1}^n p_i \le 1$ and arrows $f_i\colon A\to X_i$, there exists a unique arrow $h\colon A \to Z$ such that $h;\pi_i = p_i\cdot f_i$ for all $i=1,\ldots,n$. Observe that, by definition, the 0-ary convex product is a final object. Hereafter, we will denote the unique arrow $h$ by $\langle f_1, \dots, f_n \rangle_{\vec{p}}$ where $\vec{p}$ is a compact notation for $p_1, \dots , p_n$.

\begin{lemma}\label{lemma:naryconvexproduct}
Let $\Cat{C}$ be a $\Cat{PCA}$-enriched category with final object and binary convex products. Then, for all $n\in \mathbb{N}$, $\Cat{C}$ has $n$-ary convex products.
\end{lemma}
%




%\begin{remark}
Recall from Remark \ref{rmk:collapsed} that pcas that additionally satisfy the axiom \eqref{eq:unitstar} collapse to commutative idempotent monoids.
Observe that whenever a category $\Cat{C}$ is enriched over such pcas, it holds that $p\cdot f = f$ for all arrows $f$ and $p\in(0,1)$ and thus convex products collapse to standard products:
%\end{remark}
\begin{proposition}\label{prop:productvsconvex}
Let $\Cat{C}$ be a category enriched over commutative idempotent monoids. Let $X_1,X_2,Z$ be objects of $\Cat{C}$, and let $\pi_1\colon Z \to X_1$ and $\pi_2\colon Z \to X_2$ be arrows of $\Cat{C}$. 
\begin{center}$(Z,\pi_1,\pi_2)$ is a  product  of $X_1$ and $X_2$ iff it is a convex product of $X_1$ and $X_2$.\end{center}
\end{proposition}


By means of the proposition above, the reader can easily see that $\Cat{Rel}$, the category of sets and relations, has convex products, although these are just ordinary products. A  category with more interesting convex products is $\KlD$.

\begin{proposition}
$\KlD$ has convex products.
\end{proposition}
\begin{proof}
Let $X_1$ and $X_2$ be two objects of $\KlD$.
We show that  $X_1\oplus X_2$ with the projections $\pi_i\colon X_1 \oplus X_2 \to X_i$ for $i=1,2$ given by 
\[ \pi_1(x_1|z)\defeq \begin{cases}
    1 & z=\iota_1(x_1)\\
    0 & otherwise
\end{cases} \qquad \pi_2(x_2|z)\defeq \begin{cases}
     1 & z=\iota_2(x_2)\\
        0 & otherwise
\end{cases}\]
 is a convex product of $X_1$ and $X_2$. Above, the functions $\iota_i\colon X_i \to X_1 \oplus X_2$ denote the obvious embeddings.
 
For all arrows $f_1\colon A\to X_1$, $f_2\colon A\to X_2$ and $p_1,p_2 \in (0,1)$ such that $p_1+p_2\le 1$, the arrow $\langle f_1,f_2\rangle_{p_1,p_2}\colon A \to X_1\oplus X_2$ is defined for all $a\in A$ and $z\in X_1\oplus X_2$ as follows:
\[ \langle f_1,f_2\rangle_{p_1,p_2}(z|a) \defeq \begin{cases}
    p_1\cdot f_1(x_1|a) & z=\iota_1(x_1)\\
    p_2\cdot f_2(x_2|a) & z=\iota_2(x_2) \text{.}
\end{cases}\]
Observe that, for $i\in \{1,2\}$, $x_i\in X_i$ and $a\in A$
\begin{align*}
\langle f_1,f_2\rangle_{p_1,p_2} ; \pi_i (x_i | a) = & \sum_{z\in X_1\oplus X_2} \langle f_1,f_2\rangle_{p_1,p_2}(z|a) \cdot \pi_i (x_i | z) \\
= & \langle f_1,f_2\rangle_{p_1,p_2} (\iota_i(x_i)|a)\\
= & p_i\cdot f_i (x_i|a)
\end{align*}
namely, it holds that $(\langle f_1,f_2\rangle_{p_1,p_2} );\pi_i = p_i\cdot f_i$. In order to prove that this is the unique arrow with such a property, take $h'\colon A \to X_1\oplus X_2$ such that $h';\pi_i = p_i\cdot f_i$. For all $x_1\in X_1$ and $a\in A$, it holds that
\begin{align*}
  h'(\iota_1(x_1)|a) 
    & =\sum_{z\in X_1\oplus X_2} h'(z|a) \cdot \pi_1(x_1|z) \tag*{}\\
   &= h' ; \pi_1(x_1|a) \\
    &= p_1\cdot f_1(x_1|a) \tag*{}
\end{align*}
and similarly for $x_2\in X_2$, one obtains that $h'(\iota_2(x_2)|a) = p_2\cdot f_2(x_2|a)$, i.e., $h' = \langle f_1,f_2\rangle_{p_1,p_2}$. 
\end{proof}


\subsection{Convex biproduct categories and their monoidal algebra}\label{ssec:convexbiproduct} In both examples of categories with convex products, $\Cat{Rel}$ and $\KlD$, convex products are carried by coproducts (recall that in $\Rel$ products and coproducts coincide). In this work, we are interested in those categories where convex product and coproduct coincide, as formalised in  the following definition.  

\begin{definition}\label{def:convbicat}%\marginpar{Vedi composizione di proiezione con iniezioni}
A \emph{convex biproduct category} is a $\Cat{PCA}$-enriched category $\Cat{C}$ with an object $\zero$ which is both initial and final and, for every pair of objects $X_1,X_2$, an object $X_1\oplus X_2$ and morphisms $\pi_i \colon X_1\oplus X_2 \to X_i$ and $\iota_i \colon X_i \to X_1 \oplus X_2$ such that $(X_1\oplus X_2, \iota_1, \iota_2)$ is a coproduct, $(X_1\oplus X_2, \pi_1, \pi_2)$ is a convex product and
\begin{equation}\label{eq:delta}\iota_i; \pi_j = 
\begin{cases}  \id{X_i} & \text{if }i=j\text{,}\\
\star_{X_i,X_j} & \text{otherwise.}
\end{cases} \end{equation}
%where $\delta_{i,j} \colon X_i \to X_j$ is defined as $\delta_{i,j}\defeq \id{X_i}$ if $i=j$ and $\delta_{i,j}\defeq\star_{X_i,X_j}$ otherwise. 
A morphism of convex biproduct categories is a $\Cat{PCA}$-enriched functor $F\colon \Cat{C} \to \Cat{D}$ preserving finite coproducts. We write $\Cat{CBCat}$ for the category of convex biproduct categories and their morphisms.
\end{definition}

The above definition is obtained from that of \emph{category with finite biproducts} by just replacing products by convex products.  As for categories with finite biproducts, morphisms of convex biproduct categories preserve (convex) products. %(see Proposition \ref{prop: functor1} in Appendix~\ref{app:sec:cbproducts}). 

 \begin{proposition}\label{prop:functor 1 e mezzo}
     Let $F\colon \Cat{C}\to \Cat{D}$ be a morphism between convex biproduct categories. 
     $F$ preserves n-ary convex products.
\end{proposition}

One can easily see that both $\Cat{Rel}$ and $\KlD$ are convex biproduct categories by checking \eqref{eq:delta}. In Appendix~\ref{app:examples}, we show that the continuous analogue of $\KlD$, namely substochastic Markov kernels on standard Borel spaces, is also a convex biproduct category.
%


\medskip

Since any convex biproduct category $\Cat{C}$ has finite coproducts, it carries a symmetric monoidal category $(\Cat{C}, \oplus, \zero)$. Hereafter we will denote the structural isomorphisms by $\alpha_{X,Y,Z}\colon (X\oplus Y)\oplus Z \to X \oplus (Y \oplus Z)$, $\lambda_X\colon 0 \oplus X \to X$, $\rho_X \colon 0 \oplus X \to X$ and the symmetries by $\sigma_{X,Y}\colon X\oplus Y \to Y \oplus X$.

For all objects $X$, we define $\codiag{X}\colon X \oplus X \to X$ as the copairing of the identities $[\id{X},\id{X}]$ and $\cobang{X} \colon 0 \to X$ as the unique arrow from the initial object. Similarly, 
 for all $p\in (0,1)$, we define  $\diagp{X} \colon X \to X \oplus X$ as $\langle id_X, id_X \rangle_{p,1-p}$, and $\bang{X}\colon X \to \zero$ as the unique map to the final object $\zero$. We extend $\diagp{X}$ to arbitrary $p\in [0,1]$: $\diagpX{0\;\;}{X}\defeq \cobang{X}\oplus \id{X}$ and $\diagpX{1\;\;}{X}\defeq \id{X} \oplus \cobang{X}$.

By Fox's theorem \cite{fox1976coalgebras}, $(X, \codiag{X}, \cobang{X})$ forms a \emph{natural and coherent commutative monoid} in the monoidal category  $(\Cat{C}, \oplus, \zero)$, namely the laws in Figures~\ref{fig:monoidax} and \ref{fig:fccoherence} in Appendix~\ref{app:coherence axioms} hold. Similarly $(X,\diagp{X}, \bang{X} )$ forms a \emph{natural and coherent co-pca} in $(\Cat{C}, \oplus, \zero)$, that is the laws in Figures~\ref{fig:co-pca axioms} and \ref{fig:copcacoherence} hold. To prove the latter fact, it is convenient to first illustrate some properties of convex biproduct categories.

\begin{lemma}\label{lemma:initialproperties}
Let $\Cat{C}$ be a convex biproduct category. The following equalities hold for all objects $Z,X,Y$ in $\Cat{C}$, arrows $f,g\colon X \to Y$ and $p\in [0,1]$:
\begin{enumerate}
\begin{multicols}{2}
\item[(1)]\mylabelbis{lemma:starxy}{1} $\star_{X,Y} = \bang{X};\cobang{Y}$
\item[(3)]\mylabelbis{lemma:pi1}{3}  $\pi_1=(\id{X_1}\oplus \bang{X_2});\runit{X_1}$;
\item[(5)]\mylabelbis{lemma:iotai1}{5}  $\iota_1=\Irunit{X_1}; (\id{X_1}\oplus \cobang{X_2})$;
\item[(7)]\mylabelbis{lemma:pf}{7} $p \cdot f =\, \diagp{X}; (f \oplus \bang{X} ) ; \runit{Y} $;
\item[(2)]\mylabelbis{eq: assioma aggiuntivo}{2} $f+_p g = \langle f,g \rangle_{p,1-p};[\id{Y},\id{Y}]$;
\item[(4)]\mylabelbis{lemma:pi2}{4} $\pi_2 =(\bang{X_1}\oplus \id{X_2});\lunit{X_2}$;
\item[(6)]\mylabelbis{lemma:iota2}{6} $\iota_2 = \Ilunit{X_2} ;(\cobang{X_1}\oplus \id{X_2})$;
 \item[(8)]\mylabelbis{lemma:1-pf}{8} $(1-p)\cdot  f=\, \diagp{X} ;  \, (\bang{X}  \oplus f \,)  ; \lunit{Y}$; 
 \end{multicols}
\end{enumerate}
\end{lemma}
\begin{proposition}\label{lemma: copca objects in convbicat}
Let $\Cat{C}$ be a convex biproduct category. Then for every object $X$,  
\begin{enumerate}
 \item the triple  $(X, \codiag{X}, \cobang{X})$ is a natural and coherent monoid, i.e., the axioms in Figures~\ref{fig:monoidax} and \ref{fig:fccoherence} hold;
\item the triple $(X,\diagp{X},\bangp{X})$ is a natural and coherent co-pca, i.e., the axioms in Figures~\ref{fig:co-pca axioms} and \ref{fig:copcacoherence} hold. %\marginpar{Fare un'altra figura con naturalia' ed aggiungere qua il riferimento}
 \end{enumerate}
\end{proposition}

For all sets $X$, co-pcas and monoids in $\KlD$ are illustrated below.
\begin{equation}\label{ex:comonoids}
\arraycolsep=2pt%\def\arraystretch{2.2}
\begin{array}{cccc}
\begin{array}{rcl}
\diagp{X} \colon  X & \to & X \oplus X \\
x & \mapsto & \delta_{(x,0)}+_p\delta_{(x,1)}%\begin{cases} (x,0)\mapsto p\\ (x,1)\mapsto 1-p\end{cases}
\end{array}
&
\begin{array}{rcl}
\bang{X}  \colon X & \to & \zero\\
x & \mapsto& \star
\end{array}
&
\begin{array}{rcl}
\cobang{X}\colon  \zero &\to& X\\
\text{ }
\end{array}
&
\begin{array}{rcl}
\codiag{X}  \colon  X\oplus X & \to & X\\
(x,i) & \mapsto& \delta_{x}
\end{array}
\end{array}
\end{equation}
Monoids and co-pcas will be useful later to freely generate convex biproduct categories. In particular, morphisms of convex biproduct categories  can be characterised as follows.

\begin{proposition}\label{prop: monoidal functors}
A functor $F\colon\Cat{C}\to \Cat{D}$ is a morphism of convex biproduct categories if and only if it is a strong monoidal functor preserving monoids and co-pcas.
\end{proposition}

%We have relegated the axioms in Figures \ref{fig:monoidax}, \ref{fig:fccoherence}, \ref{fig:co-pca axioms} and \ref{fig:copcacoherence} to the Appendix, since these laws appear in Table \ref{fig:freestrictfccat} and \ref{fig:freecopcacat} for the cacse of monoidal categories that are \emph{strict} (i.e., the structural isomorphism  are identities). Indeed, hereafter, we will always tacitly assume strictness. On the one hand, this assumption is without loss of generality (by \cite{}); on the other, it allows to radically simplify calculations and the sound use of diagrammatic notations. 




\subsection{Almost a convex Fox theorem}\label{ssec:almost}
So far, we have shown that in every convex biproduct category each object carries natural and coherent structures of monoid and co-pca (Proposition~\ref{lemma: copca objects in convbicat}). Moreover, morphisms of convex biproduct categories can be characterised as precisely those monoidal functors that preserve these algebraic structures (Proposition~\ref{prop: monoidal functors}).
At this point, one might ask whether the converse holds: does any monoidal category equipped with natural and coherent monoid and co-pca structures give rise to a convex biproduct category? An affirmative answer would yield a characterisation of convex biproduct categories in the spirit of Fox's theorem~\cite{fox1976coalgebras}.
Unfortunately, this is not the case. As we shall see, an additional condition is required: the projections must be \emph{jointly monic}, i.e., for all $f,g\colon A \to X_1\oplus X_2$ the following implication holds.
\begin{center}
if $f;\pi_1=g;\pi_1$ and $f;\pi_2=g;\pi_2$ then $f=g$.
\end{center}
This condition, established in Proposition~\ref{prop:A}, will nevertheless serve as a useful principle throughout the development of our results.


We begin with the following definition, which fixes coprojections, projections, and the enrichment structure as described in Lemma~\ref{lemma:initialproperties}.
\begin{definition}\label{definition:structural}
Let $(\Cat{C}, \oplus, \zero)$ be a symmetric monoidal category such that every object $X\in\Cat{C}$ carries a natural and coherent monoid $(X,\codiag{X},\cobang{X})$ and co-pca $(X,\diagp{X},\,\bangp{X})$,  i.e., the axioms in Figures~\ref{fig:monoidax} -\ref{fig:copcacoherence} hold. 
For all objects $X_1,X_2$, $\iota_i\colon X_i \to X_1\oplus X_2$ and $\pi_i\colon X_1\oplus X_2 \to X_i$ are defined as
\begin{equation} \iota_1 \defeq \Irunit{X_1};(\id{X_1}\oplus \cobang{X_2})\text{ and } \iota_2 \defeq \Ilunit{X_2};(\,\cobang{X_1}\oplus \id{X_2})\end{equation}
\begin{equation}\label{eq:defpi}\pi_1 \defeq (\id{X_1}\oplus \bang{X_2});\runit{X_1} \quad \text{ and } \quad \pi_2 \defeq (\bang{X_1}\oplus \id{X_2});\lunit{X_2}\end{equation}
and, for all $f,g\colon X\to Y$,  $f +_p g\colon X \to Y$ and $\star_{X,Y}\colon X \to Y$ are defined as
\begin{equation}\label{eq:enrichment} f +_p g \defeq\, \diagp{X};(f \oplus g);\codiag{Y} \quad \text{ and } \quad \star_{X,Y} \defeq \bangp{X};\cobang{Y} \end{equation}
\end{definition}

Using the axioms of monoids and co-pcas (Figures~\ref{fig:monoidax}-\ref{fig:copcacoherence}), one verifies that the operations in~\eqref{eq:enrichment} endow $\Cat{C}$ with a $\Cat{PCA}$-enrichment.

\begin{lemma}\label{lemma:enrichmentpca}
Let $(\Cat{C}, \oplus, \zero)$ be a category as in Definition \ref{definition:structural}.  
Then $\Cat{C}$ is $\Cat{PCA}$-enriched. % via \ref{eq:enrichment}
\end{lemma}


Furthermore, by Fox's theorem~\cite{fox1976coalgebras}, $(\Cat{C}, \oplus, 0)$ has finite coproducts; naturality of $\bangp{X}$ implies that $0$ is also a final object. Straightforward algebraic manipulations also establish~\eqref{eq:delta}.

\begin{lemma}\label{lemma:coproducts}
Let $(\Cat{C},\oplus,0)$ be a symmetric monoidal category as in Definition~\ref{definition:structural}. %,  i.e., the axioms in Figures~\ref{fig:monoidax},\ref{fig:fccoherence} and \ref{fig:co-pca axioms},\ref{fig:copcacoherence} hold. 
Then:
\begin{enumerate} 
\item for all objects $X_1,X_2$, $(X_1\oplus X_2, \iota_1,\iota_2)$ is a coproduct;
\item $0$ is initial object;
\item $0$ is a final object;
\item the equality \eqref{eq:delta} holds.
\end{enumerate} 
\end{lemma}

Now, to conclude that a monoidal category $(\Cat{C}, \oplus, \zero)$ as in Definition \ref{lemma:enrichmentpca} is a convex biproduct category, one should only prove that $(X_1\oplus X_2,\pi_1,\pi_2)$ is a convex product. 

Given $p,q\in [0,1]$ such that $p+q\le 1$,  define $ \diaggen{p,q\,}{X}\colon X \to X\oplus X $  as  
\begin{equation}\label{diagpq generalizzata}
    \diaggen{p,q\,}{X} \defeq\,\diagp{X};(\id{X}\oplus\, \diaggen{q'}{X});(\id{X}\oplus (\id{X} \oplus\bang{X}));(\id{X}\oplus \runit{X})
\end{equation}
where $q'=\frac{q}{1-p}$ if $p\not= 1$ and $q'=0$ otherwise.

\begin{lemma}\label{lemma:diagpq}
Let $(\Cat{C}, \oplus, \zero)$ be a symmetric monoidal category as in Definition~\ref{definition:structural}. %symmetric monoidal category such that every object $X\in\Cat{C}$ carries a natural and coherent monoid $(X,\codiag{X},\cobang{X})$ and co-pca $(X,\diagp{X},\,\bangp{X})$. %,  i.e., the axioms in Figures~\ref{fig:monoidax},\ref{fig:fccoherence} and \ref{fig:co-pca axioms},\ref{fig:copcacoherence} hold. 
Then: % for $f_1\colon X\to Y_1$ and $f_2\colon X \to Y_2$:
    \begin{enumerate}
        \item\label{lemma:diagpq1} $\diaggen{p,q\,}{X};\pi_1 = p\cdot \id{X}$;
        \item\label{lemma:diagpq2} $\diaggen{p,q\,}{X};\pi_2 = q\cdot \id{X}$;
       \item\label{lemma:diagpq4} $\diaggen{p,q\,}{X};(f_1\oplus f_2); \pi_1= p\cdot f_1$ for all $f_i\colon X \to Y_i$;
       \item\label{lemma:diagpq5} $\diaggen{p,q\,}{X};(f_1\oplus f_2); \pi_2= q\cdot f_2$ for all $f_i\colon X \to Y_i$.
    \end{enumerate}
\end{lemma}

By the last two items of the lemma above, we know that there exists \emph{at least one} arrow $h$, namely $\diaggen{p,q\,}{X};(f_1\oplus f_2)$, such that $h;\pi_1=p\cdot f_1$ and $h;\pi_2=q\cdot f_2$. However, uniqueness is not guaranteed. To obtain convex products, we need to assume that the projections are jointly monic.

\begin{proposition}\label{prop:A}
Let $(\Cat{C},\oplus,0)$ be a symmetric monoidal category as in Definition~\ref{definition:structural}. 
\begin{center}If the projections $\pi_1$ and $\pi_2$ are jointly monic, then $\Cat{C}$ is a convex biproduct category.\end{center}
\end{proposition}
\begin{proof}
By Lemmas~\ref{lemma:enrichmentpca} and \ref{lemma:coproducts}, we only need to prove that $(X_1\oplus X_2, \pi_1,\pi_2)$ is a convex product. 

Let $p,q\in [0,1]$ such that $p+q\le 1$ and $f_1\colon Z\to X_1$ and $f_2\colon Z\to X_2$.
By the last two items of Lemma \ref{lemma:diagpq}, $\diaggen{p,q\,}{Z};(f_1\oplus f_2); \pi_1= p\cdot f_1$ and $\diaggen{p,q\,}{Z};(f_1\oplus f_2); \pi_2= q\cdot f_2$. To prove that this is the unique arrow with such property, suppose that there is an arrow $h\colon Z \to X_1\oplus X_2$ such that $h;\pi_1 = p\cdot f_1$ and $h;\pi_2 = q\cdot f_2$. Then, since $\pi_1$ and $\pi_2$ are jointly monic, it follows that $h =\, \diaggen{p,q\,}{Z};(f_1\oplus f_2)$.
\end{proof}
The converse does not hold: in a convex biproduct category, projections need not be jointly monic.
Finally, Lemma~\ref{lemma:diagpq} yields the following characterisation of mediating morphisms in any convex biproduct category.

\begin{lemma}\label{lemma:diagpq3} 
Let $\Cat{C}$ be a convex biproduct category. For all  $p,q\in [0,1]$ such that $p+q\le 1$ and $f_1\colon Z\to X_1$ and $f_2\colon Z\to X_2$, it holds that
 \[\langle f_1, f_2\rangle_{p,q} =\, \diaggen{p,q\,}{Z}; (f_1 \oplus f_2)\text{.}\]
\end{lemma}



\begin{remark}
The axioms in Figures~\ref{fig:monoidax}, \ref{fig:fccoherence}, \ref{fig:co-pca axioms}, and \ref{fig:copcacoherence} have been placed in the Appendix, since their \emph{strict} counterparts already appear in Tables \ref{fig:freestrictfccat} and \ref{fig:freecopcacat}. From now on, we assume monoidal categories to be strict, meaning that the structural isomorphisms  $\assoc{X}{Y}{Z}$, $\lunit{X}$, $\runit{X}$ are identities. This assumption is made without loss of generality \cite{mac_lane_categories_1978}, it greatly simplifies calculations and allows for diagrammatic notations: see the last three rows of Figure~\ref{fig:tapesax} for a diagrammatic representation of the axioms of natural monoids and co-pcas.
\end{remark}




\begin{table}[t]
\resizebox{\textwidth}{!}{%
\begin{tabular}{cc cc}
    \toprule
    $(\id{ P}\piu \codiag{ P}) ; \codiag{ P} = (\codiag{ P}\piu \id{ P}) ; \codiag{ P}$ & (\newtag{$\codiag{}$-as}{eq:codiag assoc}) &  $(\cobang{ P}\piu \id{ P}) ; \codiag{ P}  = \id{ P}  $ & (\newtag{$\codiag{}$-un}{eq:codiag unital}) \\[0.3em]
    $\cobang{\zero} = \id{\zero} \qquad\codiag{\zero} = \id{\zero}$ & (\newtag{$\cobang{\zero},\codiag{\zero}$-coh}{eq:codiag zero coherence}) & $\sigma_{ P, P};\codiag{ P}=\codiag{ P}$ & (\newtag{$\codiag{}$-sym}{eq: codiag symmetry}) \\[0.3em]
    $\cobang{P \piu Q} = \cobang{P} \piu \cobang{Q}$ & (\newtag{\,$\cobang{}$-coh}{eq:cobang coherence}) & $\codiag{P \piu Q} = (\id{P} \piu \sigma_{Q,P} \piu \id{Q}) ; (\codiag{P}\piu \codiag{Q}) $ & (\newtag{$\codiag{}$-coh}{eq:codiag coherence}) \\[0.3em]
    $\cobang{P};f =\cobang{Q}$ & (\newtag{\,$\cobang{}$-nat}{eq:cobang nat}) & $\codiag{P};f =(f\piu f); \codiag{Q}$ & (\newtag{$\codiag{}$-nat}{eq:codiag nat}) \\
    \bottomrule
\end{tabular}}
\caption{Axioms for natural and coherent monoids in a strict monoidal category.}\label{fig:freestrictfccat}
\end{table}




\begin{table}[t]
\resizebox{\textwidth}{!}{%
\begin{tabular}{cc cc}
    \toprule
    $\diagp{P};(\diagq{P}\piu \id{P})=\, \diagptilde{P};(\id{P}\piu \,\diagqtilde{})$ & (\newtag{$\diagp{}$-as}{eq:diagp assoc}) &  $\tilde{p}= pq\qquad \tilde{q}= \frac{p(1-q)}{1-pq}$ &  \\[0.3em]
    $\diagp{P};\codiag{P}  = \id{ P}$ & (\newtag{$\diagp{}$-idem}{eq:diagp idempotency}) & $\diagp{P}\sigma_{ P, P}=\,\diagpbar{P}$ & (\newtag{$\diagp{}$-sym}{eq:diagp symmetry}) \\[0.3em]
    $\diagp{\zero} = \id{\zero}$ & (\newtag{$\diagp{\zero}$-coh}{eq:diagp zero coherence}) & $\bangp{\zero} = \id{\zero}$ & (\newtag{\,$\bangp{\zero}$-coh}{eq:bangp0 coherence}) \\[0.3em]
    $\bangp{P \piu Q} = \bangp{P} \piu\, \bangp{Q}$ & (\newtag{\,$\bangp{}$-coh}{eq:bangp coherence}) & $\diagp{P \piu Q} = (\diagp{P}\piu\, \diagp{Q});(\id{P} \piu \sigma_{P,Q} \piu \id{Q}) $ & (\newtag{$\diagp{}$-coh}{eq:diagp coherence}) \\[0.3em]
    $f; \bangp{Q}=\bangp{P}$ & (\newtag{\,$\bangp{}$-nat}{eq:bangp nat}) & $f; \diagp{Q}=\,\diagp{P}; (f \piu f)$ & (\newtag{$\diagp{}$-nat}{eq:diagp nat}) \\
    \bottomrule
\end{tabular}}
\caption{Axioms for natural and coherent co-pcas in a strict monoidal category.}\label{fig:freecopcacat}
\end{table}


%The structures illustrated above provide the fundamental linguistic constructs for the language that we are going to introduce. Of course they are not arbitrary: 
%the structures in \eqref{ex:products} and \eqref{ex:comonoids} live in the Kleisli category of any \emph{symmetric monoidal monad}; those in \eqref{ex:op} are specific to the monad \(\subdistr\), but analogous structures can similarly be found in monads equipped with an algebraic presentation $\mathbb{T}$.
%
%Before introducing our diagrammatic language in Section \ref{sec:tapediagrams}, we study how these structure interact:  $\per, \uno,\copier{},\discharger{}$ provide a copy discard category (Definition \ref{def:cd}); $\piu, \zero,\codiag{},\cobang{}$  a finite coproduct category (Definition \ref{def:fp});
% together they form what we named a \emph{finite coproduct-copy discard rig category} (Definition \ref{def: distributive gs-monoidal category}). Then, by adding the arrows provided by the algebraic theory $\mathbb{T}$, one obtains a \emph{$\mathbb{T}$-copy discard rig category} (Definition \ref{def:Trig cat}).
%
%

